% Presentación final - Proyecto 5G SA (srsRAN + Open5GS + SDR), Fase 3
% Autor: Jose David Rodriguez Pinilla (A00407151) - Sistemas y Comunicaciones 2, Universidad Icesi
% Compilar con pdfLaTeX (Overleaf o texlive/texlive:latest-full), igual que los informes.
% Las cifras son las de docs/informe_fase3.tex.

\documentclass[aspectratio=169,10pt]{beamer}
\usetheme{Madrid}
\usecolortheme{default}
\usepackage[utf8]{inputenc}
\usepackage[T1]{fontenc}
\usepackage[spanish,es-noshorthands]{babel}
\usepackage{booktabs}
\usepackage{tikz}
\usetikzlibrary{arrows.meta,positioning,fit,backgrounds}
\usepackage{pgfplots}
\pgfplotsset{compat=1.16}
\setbeamertemplate{navigation symbols}{}

\newcommand{\pendiente}[1]{\textcolor{red}{\textbf{[#1]}}}

% Enlace al video y nPerf a 1 m (los mismos del informe)
\newcommand{\enlacevideo}{{\footnotesize\url{https://drive.google.com/file/d/1-HOV4RiCsWxqCv6UE2BTmOUK0HAWfsrs/view?usp=sharing}}}
\newcommand{\nperfdl}{15{,}38~Mbit/s}
\newcommand{\nperful}{7{,}82~Mbit/s}
%-------------------------------------------------------------------------------------------------

% Datos automáticos de 'make datos-fase3' (mismo archivo que usa el informe)
\IfFileExists{fase3_datos.tex}{% Generado por python -m analisis.fase3; no editar a mano (se sobrescribe).

\def\datosvigilar{Con \texttt{make vigilar} durante 44{,}4~min hubo 1 caída(s) de la B200, recuperadas solas en 6{,}4~s en promedio (MTBF de 44{,}4~min, disponibilidad $\approx 99{,}8$\,\%); el teléfono hizo 2 registros, 2 liberaciones y 0 sesiones PDU.}
}{}
\providecommand{\datosgraficamcs}{\addplot[red,only marks,mark=square*,error bars/.cd,y dir=both,y explicit]
  coordinates {(1,25) +- (0,1)};\addlegendentry{Medido 19/11}}
\providecommand{\datosvigilar}{\pendiente{vigilancia de 30--60 min}}

\title[Red 5G SA con SDR]{Red 5G Standalone con srsRAN, Open5GS y USRP B200}
\subtitle{Fase 3: integración SDR y terminal comercial}
\author[J. D. Rodríguez]{Jose David Rodríguez Pinilla}
\institute[Icesi]{Sistemas y Comunicaciones 2 --- Universidad Icesi\\Docente: Andrés Navarro Cadavid, PhD}
\date{Noviembre de 2026}

\begin{document}

\begin{frame}
\titlepage
\end{frame}

%---------------------------------------------------------------------
\begin{frame}{El proyecto en tres fases}
\begin{columns}[T]
\column{0.32\textwidth}
\textbf{Fase 1: diseño}
\begin{itemize}\small
  \item Arquitectura srsRAN--Open5GS (N1, N2, N3, N4, N6)
  \item Modelo LDPC/Polar y MCS
  \item Simulación ZMQ y presupuesto de enlace
\end{itemize}
\column{0.32\textwidth}
\textbf{Fase 2: Zero-RF}
\begin{itemize}\small
  \item Núcleo y gNB en contenedores
  \item NEA2/NIA2, trazas campo a campo
  \item RTT 31,9 ms, TCP DL 48,6 Mbit/s
\end{itemize}
\column{0.32\textwidth}
\textbf{Fase 3: al aire}
\begin{itemize}\small
  \item USRP B200, celda n78
  \item SIM programada con pySim
  \item \alert{Teléfono comercial navegando}
\end{itemize}
\end{columns}
\vspace{1em}
\begin{block}{Objetivo de la Fase 3}
Transmisión al aire y validación con un teléfono comercial real: conectar, autenticar, navegar y medir KPIs.
\end{block}
\end{frame}

%---------------------------------------------------------------------
\begin{frame}{Montaje}
\centering
\begin{tikzpicture}[font=\small,
  caja/.style={draw,rounded corners=2pt,align=center,minimum height=1cm,inner sep=4pt},
  flecha/.style={-{Latex[length=2mm]},thick}]
\node[caja,fill=blue!8] (core) {Open5GS\\AMF, SMF, UPF, \ldots};
\node[caja,fill=blue!8,left=1.4cm of core] (gnb) {gNB srsRAN\\(contenedor)};
\node[caja,fill=gray!10,below=0.4cm of core] (mongo) {MongoDB};
\begin{scope}[on background layer]
\node[draw,dashed,rounded corners,fit=(gnb)(core)(mongo),inner sep=6pt,label={above:Portátil i5-8250U, Ubuntu 24.04 nativo, Docker}] {};
\end{scope}
\node[caja,fill=orange!15,left=1.6cm of gnb] (b200) {USRP B200\\AD9364};
\node[caja,fill=green!12,below=1.4cm of b200] (tel) {OnePlus Nord\\CE 2 Lite 5G};
\node[caja,right=1.2cm of core] (inet) {Internet};
\draw[flecha,<->] (b200) -- node[above]{USB 3} (gnb);
\draw[flecha,<->] (gnb) -- node[above]{N2/N3} (core);
\draw[flecha,<->] (core) -- node[above]{N6} (inet);
\draw[flecha,<->] (b200) -- node[left,align=center]{n78 TDD\\al aire} (tel);
\draw[thick] (core) -- (mongo);
\end{tikzpicture}
\vspace{0.5em}
\begin{itemize}\small
  \item Mismos contenedores de la Fase 2: solo cambia la plantilla del gNB (ZMQ $\rightarrow$ UHD).
  \item Linux nativo: \texttt{usbipd} (WSL2) daba 18 MB/s; la celda exige $\approx$138 MB/s con \texttt{sc12}.
\end{itemize}
\end{frame}

%---------------------------------------------------------------------
\begin{frame}{Configuración de la celda}
\begin{columns}[T]
\column{0.5\textwidth}
\small
\begin{tabular}{ll}
\toprule
Banda & n78 TDD \\
Frecuencia & 3489,42 MHz \\
Ancho / SCS / PRB & 20 MHz / 30 kHz / 51 \\
Muestreo & 23,04 MS/s (sc12) \\
TDD (5 ms) & 6 DL + especial + 3 UL \\
PLMN & 001/010 (pruebas) \\
Seguridad & NEA2 + NIA2 \\
TX / RX & 50 dB / 20 dB \\
\bottomrule
\end{tabular}
\column{0.48\textwidth}
\begin{block}{¿Por qué n78?}
\small El teléfono solo hace SA en n78 en la práctica, y es la banda 5G de referencia en Colombia.
\end{block}
\begin{block}{Reloj}
\small Sin GPSDO ni generador: el desvío del TCXO ($\pm$7 kHz a 3,5 GHz) se valida porque el teléfono sincroniza, se registra y sostiene BLER 0--1\,\%.
\end{block}
\end{columns}
\end{frame}

%---------------------------------------------------------------------
\begin{frame}{Calibración de ganancias: la saturación del receptor}
\begin{columns}[T]
\column{0.5\textwidth}
\textbf{Síntoma (RX 40 dB, teléfono a 1 m)}
\begin{itemize}\small
  \item Liberación a los 8--13 s: NGAP \texttt{radio-connection-with-ue-lost}
  \item RSRP PUSCH = \texttt{ovl} (ADC desbordado)
  \item PHR $-16$ a $-31$ dB: el teléfono transmite al máximo
  \item Siguiente PUSCH: 10/10 errores, SNR $-24$ dB
\end{itemize}
\column{0.48\textwidth}
\textbf{Corrección: RX 20 dB}
\begin{itemize}\small
  \item PUSCH con SNR 28--29 dB, sin \texttt{ovl}
  \item A corta distancia el límite es el rango dinámico, no la sensibilidad
\end{itemize}
\textbf{TX}
\begin{itemize}\small
  \item 30 $\rightarrow$ 40 dB hasta ver la celda; 50 dB con margen
  \item 60 dB: la B200 cae a los 95 s
\end{itemize}
\end{columns}
\end{frame}

%---------------------------------------------------------------------
\begin{frame}{Tiempo real del anfitrión y estabilidad}
\begin{columns}[T]
\column{0.5\textwidth}
\textbf{Causas de los \textit{underflows} (20--60/min)}
\begin{itemize}\small
  \item Gobernador \texttt{powersave}
  \item \texttt{avahi-daemon} al 83\,\% de un núcleo
  \item Sondeo DRM, C-states de hasta 890 \textmu s
  \item Suspensión automática del USB
\end{itemize}
$\Rightarrow$ \texttt{make preparar-lab}: 7--9/min
\column{0.48\textwidth}
\textbf{Caídas del USB de la B200}
\begin{itemize}\small
  \item \texttt{make vigilar}: detecta, \texttt{usbreset}, recrea y espera NG Setup
  \item 19/11, solo USB: 18 min, 1 caída, \alert{16,5 s}
  \item 23/11, fuente 5 V / 4 A: 120 min, 3 caídas, MTBF \alert{40 min}, recuperación $\approx$17 s, disponibilidad $\approx$99,3\,\%
\end{itemize}
\textbf{10 MHz}: 0 \textit{underflows}, pero el teléfono no la ve $\Rightarrow$ 20 MHz
\end{columns}
\end{frame}

%---------------------------------------------------------------------
\begin{frame}{Conexión del teléfono: cada reintento, una causa}
\small
\begin{tabular}{p{4.2cm}p{5.2cm}p{4.2cm}}
\toprule
\textbf{Síntoma} & \textbf{Causa (evidencia)} & \textbf{Solución} \\
\midrule
\textit{Registration reject} 95 & SIM con MNC de 3 dígitos: SUCI 001/010 (log del AMF/SCP) & Red en 001/010 \\
Liberación a los 3 s & Plazo de sesión PDU de srsRAN & 30 s \\
Liberación a los 8--13 s & RX saturado (NGAP causa 21, \texttt{ovl}) & RX 20 dB \\
Registrado sin sesión PDU & Android en itinerancia (sin NAS tras el registro) & APN + datos en itinerancia \\
Se da de baja sin pedir datos & Wi-Fi encendido (23/11) & Apagar el Wi-Fi \\
\bottomrule
\end{tabular}
\vspace{1em}
\begin{alertblock}{Resultado (19/11/2026, 17:57:30 UTC)}
5G-AKA, sesión PDU con IP 10.45.0.2 (DNN \texttt{internet}) y \textbf{navegación por Internet} desde el teléfono.
\end{alertblock}
\end{frame}

%---------------------------------------------------------------------
\begin{frame}{Análisis de protocolos con el terminal real}
\begin{itemize}
  \item \textbf{Registro} (traza NEA0 de control): SUCI de esquema nulo, FOR=0, NSSAI SST 1, \textit{voice centric}.
  \item \textbf{5G-AKA con \textit{Synch failure}}: SQN de la SIM nueva $\neq$ UDM; AUTS y resincronización automática.
  \item \textbf{Seguridad}: Security Mode Command con NEA2 + NIA2; el NAS siguiente va cifrado.
  \item \textbf{Diagnóstico}: NGAP no va cifrado: la causa de la liberación (\texttt{radio-connection-with-ue-lost}) separó radio de NAS.
  \item \textbf{Sesión PDU} (traza del 19/11): 5G-GUTI, IP 10.45.0.2, TEID UPF \texttt{0xcc2b} / gNB \texttt{0x1}, QFI 1 con 5QI 9.
  \item \textbf{Tiempos}: registro $\rightarrow$ contexto 70~ms; Service Request $\rightarrow$ sesión en N2 240~ms; primer GTP-U a los 340~ms.
  \item \textbf{Navegación en N3}: 17,6~MB DL / 2,4~MB UL en 174~s (TLS, QUIC, DNS), pico de 16,2~Mbit/s.
\end{itemize}
\end{frame}

%---------------------------------------------------------------------
\begin{frame}{KPIs medidos con el teléfono}
\begin{columns}[T]
\column{0.55\textwidth}
\small
\begin{tabular}{lll}
\toprule
\textbf{KPI} & \textbf{Fase 3} & \textbf{Referencia} \\
\midrule
RTT medio & 30,9 ms & Fase 2: 31,9 ms \\
RTT mín./máx. & 17 / 36,9 ms & 0\,\% pérdidas \\
CQI / MCS DL & 13 / 24--26 & Previsto: 22 \\
BLER DL & 0--1\,\% & --- \\
SNR PUSCH & 28--29 dB & --- \\
MCS UL & 28 & srsUE: 5--6 \\
DL / UL nPerf 0,2 m & \textbf{43,00 / 23,95} & Cota: 46,7 \\
DL / UL nPerf 0,3 m & 38,12 / 23,47 & Cota: 43,8 \\
DL / UL 1 m & 15,38 / 7,82 & Cota: 16,4 \\
\bottomrule
\end{tabular}
\column{0.43\textwidth}
\begin{block}{Cota TS 38.306 (n78, TDD)}
\small $6\cdot\frac{948}{1024}\cdot\frac{12\cdot51}{T_s}\cdot0{,}86 \approx 81{,}9$ Mbit/s\\
$\times$ 0,657 (DL) $\approx$ \textbf{53,8 Mbit/s}
\end{block}
\small El UL llega a MCS 28: el límite de la Fase 2 era el emulador srsUE.
\end{columns}
\end{frame}

%---------------------------------------------------------------------
\begin{frame}{Presupuesto de enlace frente a lo medido}
\begin{columns}[T]
\column{0.58\textwidth}
\begin{tikzpicture}
\begin{axis}[width=\linewidth,height=5.2cm,
  xlabel={Distancia (m)}, ylabel={MCS DL},
  xmode=log, log basis x=2, xtick={0.5,1,2,4}, xticklabels={0{,}5,1,2,4},
  ymin=0, ymax=30, grid=major, legend pos=south west, legend style={font=\scriptsize},
  tick label style={font=\scriptsize}, label style={font=\scriptsize}]
\addplot[blue,thick,mark=*] coordinates {(0.5,27) (1,22) (2,18) (4,12)};
\addlegendentry{Predicción}
\datosgraficamcs
\end{axis}
\end{tikzpicture}
\column{0.4\textwidth}
\small
\begin{itemize}
  \item Espacio libre, 3489 MHz, NF 9 dB
  \item Ptx supuesta $-30$ dBm (sin generador)
  \item 19/11 a $\approx$1 m: 2--4 MCS por encima de lo previsto
  \item 23/11: MCS 11 a 1 m (en la mano), \alert{sin servicio a 2 m} (BLER 66\,\%); a 0,3 m, MCS 23--24 y, 40 min después a 0,5 m, MCS 6: pesan más la posición y el montaje que la distancia
\end{itemize}
\end{columns}
\end{frame}

%---------------------------------------------------------------------
\begin{frame}{Normativa, exposición y ética}
\begin{columns}[T]
\column{0.5\textwidth}
\textbf{Espectro}
\begin{itemize}\small
  \item MinTIC asigna, ANE planea y vigila (CNABF); n78 asignada a operadores
  \item Prueba al aire autorizada por el docente, solo en el laboratorio
  \item Potencia mínima, celda apagada al terminar
  \item PLMN de pruebas 001, nunca MCC 732
\end{itemize}
\column{0.48\textwidth}
\textbf{Exposición y privacidad}
\begin{itemize}\small
  \item B200: miliwatios; muy por debajo de ICNIRP 2020 (un teléfono emite hasta 200 mW)
  \item K y OPc solo en \texttt{.env} local, fuera de git, trazas y comandos (Ley 1581)
  \item SUCI nulo solo aceptable con SIM de laboratorio
\end{itemize}
\end{columns}
\end{frame}

%---------------------------------------------------------------------
\begin{frame}{Demostración}
\centering
\Large Video de la demostración\\[0.6em]
\normalsize \enlacevideo\\[1.2em]
\begin{itemize}
  \item Celda n78 al aire con la B200 y NG Setup
  \item Teléfono en 5G SA: registro, sesión PDU e IP 10.45.0.2
  \item Navegación y prueba de velocidad (nPerf)
  \item Métricas del gNB en tiempo real y lectura de los logs del núcleo
\end{itemize}
\end{frame}

%---------------------------------------------------------------------
\begin{frame}{Conclusiones}
\begin{enumerate}\small
  \item \textbf{Viable:} un teléfono comercial sin modificar se autenticó, obtuvo sesión PDU y navegó por una red 5G SA de software libre con un SDR de bajo costo.
  \item \textbf{Calibración por medición:} RX 40 $\rightarrow$ 20 dB por saturación; TX solo hasta tener servicio con margen.
  \item \textbf{KPIs coherentes con la teoría:} nPerf 43 Mbit/s DL y 24 Mbit/s UL (80\,\% de la cota de la celda, 92\,\% de la del MCS); RTT igual al de la Fase 2; UL MCS 28; servicio hasta 1--2 m con TX 50.
  \item \textbf{Los problemas estaban en los bordes:} SIM, Android, rango dinámico del SDR y tiempo real del anfitrión; todos diagnosticados con trazas o métricas.
  \item \textbf{Propuestas:} fuente de 6 V (la de 5 V no evitó las caídas) u otro cable, B210, anfitrión con más núcleos, GPSDO, corregir EF.AD de la SIM (volver a 001/01).
\end{enumerate}
\end{frame}

\begin{frame}
\centering
\Large ¡Gracias!\\[1em]
\normalsize Repositorio, bitácora y banco de pruebas: \texttt{github.com/JOSED235/RED5G-STANDALONE}\\[0.5em]
Video: \enlacevideo
\end{frame}

\end{document}
